\section{The actual universal patching ring and its specialization maps}
\label{sec:patched-ring-v040}

Fix the original data of \texttt{rrEx}: a prime \(\ell\), the local
list \(L\), a global object \(G_0\), a Taylor--Wiles system \(T\),
an actual \(\Lambda:\mathrm{Lim}(T)\), a complete discrete valuation
ring \(O\) with finite residue field \(k\) of characteristic \(\ell\)
and its given \(\mathbf Z_\ell\)-algebra structure, the original
\(r:\mathrm{Resid}(G_0.\mathrm{frame},O,n)\), the original
\(B,u,J\), and \(A:\mathrm{TWAut}(u,J,\mathrm{Unit},D,T)\).
Write \(q=T.q\). The assumption \(\mathrm{Adapted}(T,r)\) is needed
only where explicitly invoked below. In particular it is not an
extra hypothesis of the universal-ring construction.

This chapter constructs the complete original \(\mathrm{RRObj}\) and
\(\mathrm{RMaps}\). The determinant is the full multiplicative
homogeneous polynomial law, over every commutative scalar algebra.
No trace-only replacement and no inversion of \(n!\) is made.
The one deep published input is Chenevier's determinant
representability theorem, with its finiteness hypothesis checked in
PRG5--PRG8. The other constructions, including the relative
coefficient extension and the specified local factor of the ordered
splitting algebra, are proved below.

% BEGIN PRG1
\paragraph{PRG1. The precise coefficient category and tuple to be represented.}
\label{claim:PRG1}
Let \(\mathcal C_O\) be the category whose objects are complete
Noetherian local \(O\)-algebras \(R\) such that \(O\to R\) is local
and induces an isomorphism \(k\simeq R/\mathfrak m_R\), and whose
arrows are all \(O\)-algebra homomorphisms. This is exactly the
source's \(\mathrm{IsCO}(O,R)\): its residue-surjectivity field gives
surjectivity of \(k\to R/\mathfrak m_R\), while localness gives
injectivity. Conversely the residue isomorphism says that every
\(x\in R\) differs from an element of \(O\) by an element of
\(\mathfrak m_R\), which is the specified field of \(\mathrm{IsCO}\).

Set
\[
 P=\ker(\Lambda.G.\eta\circ\Lambda.\pi_o)\le\Lambda.\mathrm{Ph},
 \quad p=\Lambda.h.p_0\circ\pi_{o,0}:P\longrightarrow G_0.G_0.
 \tag{PRG.1}
\]
Here the last \(G_0.G_0\) denotes the plus subgroup of the original
global object; \(\pi_{o,0}\) is the source's codomain-restricted
\(\pi_o\), not a newly chosen map. A tuple over \(R\) consists of
\[
 (\mathcal D,(\psi_{v,i}),s,(U_{j,i}),(y_{j,i})),                 \tag{PRG.2}
\]
where \(\mathcal D\) is a continuous rank-\(n\) determinant on
\(P\) reducing to \(r.\overline D\circ p\), each \(\psi_{v,i}\)
satisfies the whole original \(r.\mathrm{IsPsi}\), and
\[
 s,y_{j,i},U_{j,i}-A.\mathrm{ut}_{j,i}\in\mathfrak m_R.
 \tag{PRG.3}
\]
Constants in (PRG.3) are mapped from \(O\). A morphism of tuples
must carry all five entries of (PRG.2) to their counterparts; equality
of the first entry is equality of full laws. These are exactly
\texttt{Tup}, \texttt{Tup.Maps}, and the universal property in
\texttt{RRObj}. The coefficient category does not impose an
unmentioned continuity restriction on its arrows: PRG2 proves it
from the stated fields.
% END PRG1

% BEGIN PRG2
\paragraph{PRG2. Automatic locality, residue compatibility and continuity.}
\label{claim:PRG2}
Every \(O\)-algebra map \(f:R\to R'\) between objects of
\(\mathcal C_O\) is local, induces the identity on the identified
residue field \(k\), and is continuous for the maximal-ideal
topologies. To prove the nontrivial assertion, let
\(x\in\mathfrak m_R\) and choose \(b\in O\) with
\(f(x)-b\in\mathfrak m_{R'}\). If \(b\) were a unit of \(O\),
then \(x-b\) would be a unit of \(R\), whose image is a nonunit of
\(R'\), impossible. Thus \(b\in\mathfrak m_O\), and localness of
\(O\to R'\) gives \(f(x)\in\mathfrak m_{R'}\). Units map to
units, so \(f^{-1}\mathfrak m_{R'}=\mathfrak m_R\).
Subtracting an \(O\)-lift of the residue of any element now proves
residue compatibility. Multiplication gives
\(f(\mathfrak m_R^a)\subseteq\mathfrak m_{R'}^a\) for every \(a\),
which proves continuity and the source's \texttt{res\_map} identity.

Every \(R/\mathfrak m_R^a\) is finite. Indeed the successive quotients
\(\mathfrak m_R^b/\mathfrak m_R^{b+1}\) are finite-dimensional over
the finite field \(k\), since \(\mathfrak m_R\) is finitely
generated. Induction on \(a\) gives finiteness. Completeness and
separation identify \(R\) with the inverse limit of these finite
local rings, with surjective transition maps. Thus each object of
\(\mathcal C_O\) is a profinite local ring of exactly the kind used
in PRG5. Finally a characteristic-coefficient map on a profinite
group is continuous into \(R\) exactly when all its reductions
modulo \(\mathfrak m_R^a\) are locally constant. This identifies
the source's \(\mathrm{IsCont}\) and \(\mathrm{AdicCont}\) with
the continuity in the published theorem.
% END PRG2

% BEGIN PRG3
\paragraph{PRG3. Finite power series and complete presentations.}
\label{claim:PRG3}
For a finite set \(E\), \(R[[X_e:e\in E]]\) belongs to
\(\mathcal C_O\). Its maximal ideal is
\((\mathfrak m_R,(X_e))\); the residue field is \(k\).
The coefficientwise description proves separation and completeness,
and the usual leading-coefficient induction proves the formal
Hilbert basis theorem, hence Noetherianity. If \(f:R\to R'\) is an
\(O\)-algebra map and \(a_e\in\mathfrak m_{R'}\), there is a
unique \(O\)-algebra homomorphism
\[
 R[[X_e]]\longrightarrow R',\qquad
 \sum_\nu b_\nu X^\nu\longmapsto
 \lim_{N}\sum_{|\nu|<N} f(b_\nu)a^\nu.                    \tag{PRG.4}
\]
The tails lie in \(\mathfrak m_{R'}^N\); addition and multiplication
can therefore be verified in each finite quotient. The restriction
and variable values determine every finite truncation of any other
map by PRG2, and separation gives uniqueness. Iterating this argument
also treats the original nested power series rings.

Every \(R\in\mathcal C_O\) has a presentation
\(O[[X_1,\ldots,X_d]]/I\). Choose finitely many lifts of generators
of \(\mathfrak m_R/(\mathfrak m_OR+\mathfrak m_R^2)\).
Nakayama gives \(\mathfrak m_R=(\mathfrak m_O,a_1,\ldots,a_d)R\).
Successively lifting the residue, then the degree-one error, then
each higher-degree error expresses an arbitrary element of \(R\)
as the convergent value of a power series with coefficients in \(O\).
Equivalently, at stage \(N\) correct the error in
\(\mathfrak m_R^N/\mathfrak m_R^{N+1}\) by monomials of degree
\(N\) in the chosen generators and a uniformizer of \(O\).
The corrections have orders tending to infinity in the source, so
they sum there. This proves surjectivity of (PRG.4).

We use the standard complete Noetherian module facts in this
argument and below: a finite module over a complete Noetherian local
ring is complete and separated, every submodule of a finite module
is closed, and quotients by proper ideals are complete Noetherian
local rings. Their hypotheses here are always a finite module and
the maximal-ideal topology. One proof applies Artin--Rees to a
finite presentation, making the induced and intrinsic submodule
topologies equivalent, and then takes the inverse limits of its
finite quotient presentations. Krull intersection gives separation;
the surjective transition maps give exactness at the quotient.
These facts do not require flatness of a coefficient extension.
For precise published statements see the Stacks Project,
Lemma 10.31.2 (Tag 0306), \url{https://stacks.math.columbia.edu/tag/0306},
and Lemma 10.97.1 (Tag 00MA),
\url{https://stacks.math.columbia.edu/tag/00MA}.
In the first application the base ring is Noetherian and the
variable set is finite; in the second, the base is Noetherian,
the ideal is its maximal ideal, and all modules are finite.
% END PRG3

% BEGIN PRG4
\paragraph{PRG4. The Witt coefficient ring maps to the given \(O\).}
\label{claim:PRG4}
There is a local continuous map \(W(k)\to O\) inducing the identity
on \(k\), and compatible with the given \(\mathbf Z_\ell\)-map.
Here it need not be injective. This includes equal characteristic.
For completeness choose a generator \(\alpha\) of the finite
extension \(k/\mathbf F_\ell\), its monic separable minimal
polynomial \(\overline P\), and a monic lift
\(P\in\mathbf Z_\ell[X]\). The ring
\(U=\mathbf Z_\ell[X]/(P)\) is finite free, complete, has residue
field \(k\), and maximal ideal \((\ell)\). Its being local follows
by reducing a finite integral algebra modulo \(\ell\); an element
with nonzero residue is invertible by a convergent geometric series.
Every nonzero element has a finite \(\ell\)-adic order, by its
free coefficient expansion. Dividing by that power of \(\ell\)
leaves a unit. Thus \(U\) is the unramified complete discrete
valuation ring with residue \(k\), canonically identifiable with
\(W(k)\) after the residue identification. More explicitly Hensel
lifting \(\alpha\) in \(W(k)\) gives a map \(U\to W(k)\);
both sides are free of degree \([k:\mathbf F_\ell]\) over
\(\mathbf Z_\ell\), and its residue matrix is invertible, so it
is an isomorphism.

In the given \(O\), start from any lift of \(\alpha\).
The derivative of \(P\) at that lift is a unit. Newton iteration
\(b\mapsto b-P(b)/P'(b)\) squares the order of its error and
converges in \(O\) to a root reducing to \(\alpha\). The same
derivative calculation proves uniqueness with that residue. It
defines the required map \(U\to O\). The given
\(\mathbf Z_\ell\to O\) is continuous: \(\ell\) belongs to
\(\mathfrak m_O\), so its \(a\)-th power maps into
\(\mathfrak m_O^a\). The resulting map is local by its residue
isomorphism and continuous for the same reason. We have not assumed
that \(O\) is an unramified extension of \(\mathbf Z_\ell\).
% END PRG4

% BEGIN PRG5
\paragraph{PRG5. The exact published determinant theorem used.}
\label{claim:PRG5}
We apply Chenevier, \emph{The p-adic analytic space of pseudocharacters
of a profinite group and pseudorepresentations over arbitrary rings},
Lemma 3.2, Proposition 3.3, Remark 3.5 and Proposition 3.7,
pp.~41--43 of the 56-page version at
\url{https://maths.dur.ac.uk/events/Meetings/LMS/2011/GRAF11/papers/chenevier.pdf}.
For a profinite group \(H\), a finite field \(k\) of characteristic
\(\ell\), and a continuous residual determinant of positive degree
coming from a semisimple residual representation, they provide a
universal continuous determinant over a profinite local
\(W(k)\)-algebra \(R_W\) with residue \(k\). If every open subgroup
of \(H\) has finitely many continuous homomorphisms to
\(\mathbf F_\ell\), then \(R_W\) is Noetherian and has a finite
power-series presentation over \(W(k)\).

The source's \(\mathrm{PhiCond}(H,\ell)\) is exactly this condition:
the additive group \(\mathbf F_\ell\) is written there as
\(\mathrm{Multiplicative}(\mathbf Z/\ell)\), and a homomorphism to
a finite discrete group is continuous if and only if its kernel is
open. We use the first alternative of Proposition 3.7, its
condition (F). Its other alternative, which assumes \(\ell>n\),
is not used. No assumption that \(\ell>n\), that the residual
representation is irreducible, or that \(n!\) is a unit enters this
application. The polynomial laws in the theorem are precisely full
determinants. Lemma 3.2 passes from finite local rings to their
inverse limits, so the universal property holds for all objects
in PRG2, not only for Artinian targets. PRG7 verifies (F) and the
residual representation for the actual global group; PRG10 verifies
them for every actual local group. This paragraph records a
published-theorem application, rather than claiming a new proof
of its deep Noetherianity assertion.
% END PRG5

% BEGIN PRG6
\paragraph{PRG6. Relative determinant representability over the original \(O\).}
\label{claim:PRG6}
Under the group and residual hypotheses of PRG5, continuous lifts of
the residual determinant to objects of \(\mathcal C_O\) have a
universal object in that same category. Choose a presentation
\(R_W=W(k)[[Z_1,\ldots,Z_e]]/I\) and the map in PRG4, and put
\[
 R_O=O[[Z_1,\ldots,Z_e]]/(I\,O[[Z_1,\ldots,Z_e]]).            \tag{PRG.5}
\]
The image of \(I\) is contained in the maximal ideal, because the
original presentation has residue \(k\) and PRG4 is the identity
on that residue. Thus (PRG.5) is a nonzero object of
\(\mathcal C_O\) by PRG3. Base-changing the universal law gives
a continuous law over \(R_O\) with the specified residual law.

For a target \(R'\in\mathcal C_O\), a lift gives a unique
continuous \(W(k)\)-map \(R_W\to R'\) by PRG5. Together with
the prescribed \(O\)-map it evaluates the variables in (PRG.5)
in \(\mathfrak m_{R'}\) and kills the image of \(I\). This gives
an \(O\)-map \(R_O\to R'\) carrying the universal full law to
the target. Conversely any such \(O\)-map restricts to the unique
\(W(k)\)-map; it therefore has the same values on all the \(Z_i\),
and PRG3 proves uniqueness. PRG2 shows that uniqueness is among
\emph{all} \(O\)-algebra maps, as demanded by \texttt{IsUnivDet}.
No flatness of \(W(k)\to O\) or \(O\to R'\) was used. If the
degree is zero, a unital multiplicative degree-zero law is the
constant-one law over every scalar algebra; \(O\) itself represents
the singleton functor. This separately covers \(n=0\).
% END PRG6

% BEGIN PRG7
\paragraph{PRG7. The finiteness and residual hypotheses on the actual cover.}
\label{claim:PRG7}
The group \(P\) in (PRG.1) is profinite and topologically finitely
generated. Indeed it is the kernel of the actual continuous map to
the finite sign group, hence is open and closed in
\(\Lambda.\mathrm{Ph}\). The latter is profinite and has the
specified finite topological generating set in \(\Lambda.\mathrm{fg}\).
An open subgroup of a topologically finitely generated profinite
group is topologically finitely generated: take the dense abstract
subgroup generated by the given finite set. Its intersection with
the open subgroup is dense in that subgroup and has finite index
in the abstract subgroup. Coset representatives give the finite
Schreier generating set \(tgs^{-1}\), with \(s\) the representative
of the coset of \(tg\). Rewriting a word one letter at a time proves
the generating assertion. Taking closures proves the topological
version.

This also proves \(\mathrm{PhiCond}(P,\ell)\). Every open subgroup
of \(P\) has a finite topological generating set. A continuous
homomorphism from it to a finite group is determined by its values
on that set: two such maps agree on the generated dense subgroup,
and their equality locus is closed. There are only finitely many
assignments of values. In particular there are finitely many maps
to \(\mathbf F_\ell\).

The map \(p\) is the original continuous surjection. Restricting
\(\pi_o\) to the sign kernels is surjective because its sign is
the pullback of the target sign; the same argument applies to
\(\Lambda.h.p_0\). The characters
\(\overline\chi_i=r.\mathrm{chib}_i\circ p\) have open kernels
by continuity and are distinct by surjectivity of \(p\).
Their diagonal sum is a continuous semisimple representation over
\(k\), and its determinant is exactly
\(r.\overline D\circ p\) over all scalar algebras. This verifies
every group and residual hypothesis of PRG5 for (PRG.1). Neither
\(\mathrm{Adapted}\) nor a new choice of limit is used.
% END PRG7

% BEGIN PRG8
\paragraph{PRG8. The actual global determinant deformation ring.}
\label{claim:PRG8}
Apply PRG6 and PRG7 to obtain an object \(R_{\mathcal D}\) of
\(\mathcal C_O\) and a continuous determinant
\(\mathcal D_{\mathrm{univ}}\) on the actual group \(P\).
For every target \(R'\in\mathcal C_O\) and every continuous full
determinant \(\mathcal D'\) reducing to \(r.\overline D\circ p\),
there is a unique \(O\)-algebra map
\[
 R_{\mathcal D}\longrightarrow R',\qquad
 \mathcal D_{\mathrm{univ}}\longmapsto\mathcal D'.            \tag{PRG.6}
\]
This is the original \texttt{univDet\_mf} endpoint required by
\texttt{pre\_univ}, restricted to the actual coefficient DVR
and cover of \texttt{rrEx}. The residual multiplicity-free field
has been checked in PRG7, although the version of the published
theorem used here is more general. Rank zero uses the last sentence
of PRG6. In particular the construction remains valid for the
selected rank \(4\) and prime \(3\); it does not pass through a
trace deformation functor which would require a factorial to be
invertible. The object is produced, not assumed as an input to the
remaining patching argument.
% END PRG8

% BEGIN PRG9
\paragraph{PRG9. Degree-one laws and compatibility on the exact subgroup \(\mu_v\).}
\label{claim:PRG9}
A unital multiplicative degree-one polynomial law on \(R[H]\)
is an \(R\)-algebra homomorphism. Indeed degree-one laws are
linear: expanding on independent variables leaves only linear
terms, and naturality gives additivity and scalar compatibility.
Multiplicativity and the unit condition then give an algebra map.
Its values on group elements are units and form a character
\(H\to R^\times\). Conversely the linear extension of a character
is such a law. These operations commute with arbitrary coefficient
maps; continuity of the law is exactly continuity of the character.

For the original local group \(L.G_v\), let \(\mu_v\) be the
inverse image of the \(\ell\)-primary torsion subgroup of its
topological abelianization, exactly as in \texttt{GFrame.mu}.
If \(g\in\mu_v\), some \(g^{\ell^a}\) lies in the closed
commutator subgroup. The open-kernel residual character
\(r.\overline\psi_{v,i}\) kills that subgroup: it kills
commutators and its kernel is closed. The given
\(r.\theta_{v,i}\) also kills it by \(r.\mathrm{th\_comm}\).
Consequently both \(r.\overline\psi_{v,i}(g)\) and the residue
of \(r.\theta_{v,i}(g)\) have \(\ell\)-power order in \(k^\times\).
They equal \(1\), because
\((x-1)^{\ell^a}=x^{\ell^a}-1\) in characteristic \(\ell\)
and a field has no nonzero nilpotents. Thus
\[
 r.\overline\psi_{v,i}(g)=
 \overline{r.\theta_{v,i}(g)}\quad(g\in\mu_v).              \tag{PRG.7}
\]
This verifies consistency of the exact local-character condition;
it does not replace \(\mu_v\) by inertia or prescribe an extra
extension of \(\theta\) to the whole local group.
% END PRG9

% BEGIN PRG10
\paragraph{PRG10. Every constrained local-character ring is constructed.}
\label{claim:PRG10}
For every original \(v\) and \(i\), the group \(L.G_v\) is
profinite, \(T.\mathrm{phiV}(v)\) is its full
\(\mathrm{PhiCond}\), and \(r.\overline\psi_{v,i}\) has open
kernel. PRG5 and PRG6 in degree one, followed by PRG9, give a
universal character \(\psi^0_{v,i}\) over an object
\(R^0_{v,i}\in\mathcal C_O\). Define
\[
 R_{v,i}=R^0_{v,i}/
 \bigl(\psi^0_{v,i}(g)-r.\theta_{v,i}(g):g\in\mu_v\bigr).
 \tag{PRG.8}
\]
Every displayed generator has zero residue by (PRG.7), so this is
a proper ideal. Although the displayed family can be infinite,
the ideal is finitely generated and closed by Noetherianity and
PRG3. Hence \(R_{v,i}\) belongs to \(\mathcal C_O\).
The pushed-forward character is continuous by PRG2, has the
required residue, and has precisely the prescribed values on
\(\mu_v\). It therefore satisfies all three fields of
\(r.\mathrm{IsPsi}\).

For any target character satisfying those three fields, the unique
map from \(R^0_{v,i}\) kills the ideal in (PRG.8), and so descends
uniquely. Conversely a map from (PRG.8) preserves the three fields
by PRG2 and the defining relations. This proves the complete
universal property of
\texttt{universalDeformationRing\_localCharacters} for the present
data. No absolute irreducibility theorem in higher rank is used;
degree one was identified with characters directly in PRG9.
% END PRG10

% BEGIN PRG11
\paragraph{PRG11. Finite completed coproducts in the original category.}
\label{claim:PRG11}
For a finite family \((R_a)\) in \(\mathcal C_O\), there is an
object \(R_*\) with \(O\)-maps \(j_a:R_a\to R_*\) such that
\[
 \operatorname{Hom}_O(R_*,R')
 \simeq\prod_a\operatorname{Hom}_O(R_a,R')                 \tag{PRG.9}
\]
for every \(R'\in\mathcal C_O\), with the indicated maps given
by restriction. Write \(R_a=O[[Z_a]]/I_a\) by PRG3, using disjoint
finite variable sets. Set
\[
 R_*=O[[\coprod_a Z_a]]/(I_a:\text{all }a).                 \tag{PRG.10}
\]
Each ideal maps into the maximal ideal of the common power series
ring, so the quotient is nonzero and has residue \(k\). It is
complete and Noetherian by PRG3. The canonical evaluations give
the \(j_a\). A family of maps to \(R'\) specifies all the variable
values in \(\mathfrak m_{R'}\) by PRG2; (PRG.4) evaluates them
simultaneously and kills all \(I_a\). Conversely an \(O\)-map
out of (PRG.10) is determined by these values. This proves (PRG.9)
without an unproved existence assertion about completed tensor
products. The empty family gives \(O\).
% END PRG11

% BEGIN PRG12
\paragraph{PRG12. The joint determinant and ordered local-character object.}
\label{claim:PRG12}
Apply PRG11 to \(R_{\mathcal D}\) and all \(R_{v,i}\) from
PRG8 and PRG10. There are finitely many: the local list and
\(\mathrm{Fin}(n)\) are finite. Denote the resulting object by
\(R_1\), and push the universal determinant and characters to it.
The result is a continuous full determinant \(\mathcal D_1\)
with the required residue and an ordered family \(\psi^1_{v,i}\)
satisfying \(r.\mathrm{IsPsi}\). The conditions survive the push
by PRG2. For any target determinant and ordered family satisfying
the same conditions, PRG8 and PRG10 give unique maps from their
respective rings. PRG11 joins them into one map from \(R_1\).
If another map carries the determinant and every character to
those targets, the individual universal properties force all its
restrictions, and (PRG.9) forces the map itself. Thus
\(R_1\) represents the entire joint datum, exactly the
\texttt{pre\_univ} statement used by \texttt{rrEx}.
There are no relations between the global determinant and these
local characters at this stage: the later six families of
relations impose the required global conditions. Imposing them
here would change the tuple functor and its universal property.
% END PRG12

% BEGIN PRG13
\paragraph{PRG13. Construction of the complete original \(\mathrm{RRObj}\).}
\label{claim:PRG13}
Form the finite power series ring
\[
 R=R_1[[S,(V_{j,i})_{j,i},(Y_{j,i})_{j,i}]],\qquad
 s=S,\quad U_{j,i}=A.\mathrm{ut}_{j,i}+V_{j,i},\quad
 y_{j,i}=Y_{j,i}.                                          \tag{PRG.11}
\]
Push \(\mathcal D_1\) and \(\psi^1\) along \(R_1\to R\).
The result satisfies all fields of (PRG.2)--(PRG.3) by PRG2 and
PRG3. Given a target tuple, PRG12 gives its unique map
\(f_1:R_1\to R'\). The remaining required variable values are
\[
 S\mapsto s',\qquad V_{j,i}\mapsto U'_{j,i}-A.\mathrm{ut}_{j,i},
 \qquad Y_{j,i}\mapsto y'_{j,i},                            \tag{PRG.12}
\]
which lie in \(\mathfrak m_{R'}\) by the target tuple fields.
PRG3 extends \(f_1\) uniquely to an \(O\)-map \(R\to R'\).
It carries all five entries to the target. Conversely any map
carrying all five entries must restrict to \(f_1\) by PRG12,
must have the variable values (PRG.12), and hence is the same map
by PRG3. This proves \texttt{rrEx}, with every original field of
\(\mathrm{RRObj}\), for every \(n\), every given \(\Lambda\),
and every given \(A\). In particular there is no hidden
assumption that the desired universal ring already exists.
% END PRG13

% BEGIN PRG14
\paragraph{PRG14. Independence of the presentation, and arbitrary supplied \(W\).}
\label{claim:PRG14}
Two objects satisfying the complete \(\mathrm{RRObj}\) universal
property are uniquely isomorphic by an \(O\)-algebra isomorphism
preserving the tuple. To prove this, apply each universal property
to the other object's tuple. Both composites preserve the original
tuple, so uniqueness makes them identity maps. Any other
tuple-preserving isomorphism equals these maps by the same
uniqueness. Thus the choices of presentations, polynomial lifts
and generators in PRG4--PRG13 do not alter the represented datum.

In the rest of this chapter \(W\) may be \emph{any} original
\(\mathrm{RRObj}(\Lambda,A)\), as in \texttt{rmapsEx}; it need
not be the explicit presentation (PRG.11). All constructions below
use only its actual \(O\)-algebra \(W.R\), its actual tuple
\(W.t\), and its full universal property. The isomorphism just
proved explains why this scope is stronger than a calculation for
one convenient presentation. It also prevents a choice of a new
ring from changing the finite-level or specialization targets.
% END PRG14

% BEGIN PRG15
\paragraph{PRG15. The ordered polynomial splitting algebra is finite free.}
\label{claim:PRG15}
Let \(C\) be a commutative ring and \((P_j)_{j<q}\) monic
polynomials of degree \(n\). There is a universal algebra \(E\)
with ordered roots \(r_{j,1},\ldots,r_{j,n}\) satisfying
\(P_j(X)=\prod_i(X-r_{j,i})\). It is generated as a
\(C\)-algebra by these roots and is finite free of rank
\((n!)^q\). Here is a direct construction. Adjoin the first root
by \(C[X_1]/(P_j(X_1))\), a free module with basis
\(1,X_1,\ldots,X_1^{n-1}\). Divide \(P_j(X)\) by
\(X-X_1\) in this algebra. The remainder is \(P_j(X_1)=0\),
so the quotient is monic of degree \(n-1\). Repeat, adjoining
each next root of the remaining monic quotient. The final ranks
multiply to \(n!\). Do this for each \(j\); finite free ranks
multiply again.

An ordered splitting in any target defines each successive map
uniquely. At each step the remaining factor is the product of the
remaining linear factors: multiplication by a monic polynomial
is injective in a polynomial ring over any coefficient ring, so
cancellation is valid even with zero divisors. This proves the
universal property and its uniqueness. For \(n=0\), a monic
degree-zero polynomial is \(1\), the construction adjoins no
root and has rank one. The construction makes no separability or
distinct-root assumption.
% END PRG15

% BEGIN PRG16
\paragraph{PRG16. A specified residue splitting selects a finite free local factor.}
\label{claim:PRG16}
Now suppose \(C\) is complete Noetherian local with residue field
\(k\), and that \(P_j\bmod\mathfrak m_C=\prod_i(X-\bar u_{j,i})\)
with \(u_{j,i}\in C\). The ordered roots give a surjective
\(C\)-map \(\chi:E\to k\) sending \(r_{j,i}\) to
\(\bar u_{j,i}\). Set \(\mathfrak n=\ker\chi\). Then
\(E_{\mathfrak n}\) is a nonzero finite free local
\(C\)-algebra, complete for its maximal ideal, with residue \(k\).
Moreover \(E\to E_{\mathfrak n}\) is surjective.

Indeed \(E\) is finite free and hence complete for
\(\mathfrak m_C E\) by PRG3. The finite-dimensional algebra
\(E/\mathfrak m_CE\) is Artinian and is the product of its
finitely many local factors: powers of its distinct maximal
ideals are comaximal and their intersection is zero, giving this
decomposition by the Chinese remainder theorem. Its idempotents
lift uniquely to \(E\). To see this explicitly, if \(e^2-e\)
lies in the current error ideal, \(2e-1\) is a unit modulo that
ideal and hence in \(E\), since \(\mathfrak m_C E\) is in the
Jacobson radical of a finite \(C\)-algebra. Replacing
\(e\) by \(e-(e^2-e)(2e-1)^{-1}\) squares the error.
Iteration and completeness give an idempotent. Two lifts differ
by an element of \(\mathfrak m_CE\); their standard idempotent
decompositions, or the same unit calculation, prove uniqueness.
Orthogonality and the sum-one identity therefore lift as well.

Let \(e\) select the factor singled out by \(\chi\).
Then \(eE\) is a direct summand of the finite free \(C\)-module
\(E\), hence finite projective and therefore free over the local
ring \(C\). It is local because every maximal ideal of a finite
integral \(C\)-algebra lies over \(\mathfrak m_C\), and its
residue algebra has only the selected maximal ideal. It is the
localization \(E_{\mathfrak n}\): the other idempotents die there,
while every element outside \(\mathfrak n\) is a unit in \(eE\).
The localization map is thus the surjective projection \(E\to eE\).
If \(\mathfrak m_{eE}\) is its maximal ideal, its reduction
modulo \(\mathfrak m_C\) is nilpotent; hence
\(\mathfrak m_{eE}^{a}\subseteq\mathfrak m_CeE\) for some \(a\).
The two ideal topologies are equivalent, proving maximal-ideal
completeness. Its residue is the specified \(k\), and its free
rank is positive since its residue is nonzero.
% END PRG16

% BEGIN PRG17
\paragraph{PRG17. Identification with the source's exact completed splitting algebra.}
\label{claim:PRG17}
For the data of PRG16 let
\[
 \operatorname{Sp}(P,u)=
 C[[(Z_{j,i})]]/
 \bigl(\text{coefficients of }
      \prod_i(X-u_{j,i}-Z_{j,i})-P_j(X):j<q\bigr).          \tag{PRG.13}
\]
This is precisely the source's \(\mathrm{Sp}\), with roots
\(u_{j,i}+Z_{j,i}\). Reduction at all \(Z=0\) followed by
\(C\to k\) kills the relations, so the ideal is proper and
(PRG.13) is complete Noetherian local with residue \(k\).
In the factor \(E_{\mathfrak n}\) of PRG16,
\(r_{j,i}-u_{j,i}\) lies in its maximal ideal. PRG3 evaluates
the power series there and kills the relations, giving
\(\operatorname{Sp}(P,u)\to E_{\mathfrak n}\).

Conversely the actual roots of (PRG.13) give a map
\(E\to\operatorname{Sp}(P,u)\) by PRG15. Its composite with
the residue map is \(\chi\); consequently every element of
\(E\setminus\mathfrak n\) maps to a unit. It extends to
\(E_{\mathfrak n}\to\operatorname{Sp}(P,u)\). The composites
are identities: on \(E_{\mathfrak n}\), check the roots and
the surjective projection of PRG16; on (PRG.13), check coefficients
and each \(Z_{j,i}\) and apply the power-series uniqueness in
PRG3 before taking the quotient. Thus (PRG.13) is the very same
local factor, and is finite free of positive rank over \(C\).
It is generated as an ordinary \(C\)-algebra by its roots,
because \(E\to E_{\mathfrak n}\) is surjective. If
\(C\in\mathcal C_O\), its residue field and coefficient maps
make (PRG.13) an object of \(\mathcal C_O\) as well.
This proves all four fields of \texttt{spCo}; no mere topological
generation statement has been substituted for the required
ordinary algebra generation.
% END PRG17

% BEGIN PRG18
\paragraph{PRG18. The actual \(\mathrm{Spo}(A)\) and every finite-level splitting target.}
\label{claim:PRG18}
For a full determinant of degree \(n\), the characteristic
polynomial of a group element is monic of degree \(n\), and
coefficient base change gives its residual characteristic
polynomial. These follow by evaluating the law at \(X-g\):
the leading coefficient is its value at \(1\), and the law is
homogeneous of degree \(n\). For the original pair \(u\), the
residual law is the determinant of the diagonal characters.
Thus for \(g=T.\mathrm{sg}_j\) it reduces to
\[
 \prod_i(X-r.\mathrm{chib}_i(g))
   =\prod_i(X-\overline{A.\mathrm{ut}_{j,i}}),               \tag{PRG.14}
\]
by \(A.\mathrm{ut\_res}\). PRG17 with these polynomials over
the original \(B\) gives the literal
\(\mathrm{Spo}(A)=\mathrm{SpAt}(u,T.\mathrm{sg},A.\mathrm{ut})\).
It is a nonzero object of \(\mathcal C_O\), finite free over
\(B\), and generated as a \(B\)-algebra by its actual ordered
roots.

For the second family in \texttt{spoCo}, replace \(g\) by
\((T.h_N).p_0(T.\mathrm{ph}_{N,j})\). The quotient of that element
by \(T.\mathrm{sg}_j\) lies in the original \(T.H\) by
\(T.\mathrm{ph\_H}\). Under \(\mathrm{Adapted}(T,r)\), every
\(r.\mathrm{chib}_i\) is trivial there. Hence (PRG.14) remains
valid with exactly the same \(A.\mathrm{ut}\), and PRG17 gives
all finite-level \(\mathrm{SpAt}\) targets with all four fields.
The first, limiting target did not itself need adaptedness; it is
the simultaneous finite-level assertion that uses it.
% END PRG18

% BEGIN PRG19
\paragraph{PRG19. Every actual finite-level datum is a tuple over its original ring.}
\label{claim:PRG19}
Assume \(\mathrm{Adapted}(T,r)\). Write
\(R_N=(A.A_N).R\), \(u_N=(A.A_N).p\), and
\(s_N=(A.A_N).s\). The tuple over \(R_N\) is exactly
\[
 \left(u_N.D\circ\pi_{N,0},\ u_N.\psi,\
 s_N(CX_{\mathrm{unit}}),\ (A.A_N).U,\
 (s_N(Y_{j,i}))_{j,i}\right).                             \tag{PRG.15}
\]
The map \(\pi_{N,0}\) is the original \(\pi_N\) restricted
and codomain-restricted to the plus subgroups. It is defined there
because \(\Lambda.\mathrm{modH}\), \(T.H_\eta\), and the
sign compatibilities show its sign is one. It is continuous, so
the full determinant in (PRG.15) is continuous. For every
\(g\in P\), the two elements
\((T.h_N).p_0(\pi_{N,0}g)\) and \(p(g)\) differ by an element
of \(T.H\). Adaptedness makes their values under every residual
character equal. The actual residual compatibility of \(u_N\)
then makes its pulled-back residual law exactly
\(r.\overline D\circ p\). This equality holds for the diagonal
representations and hence for their full laws, not only traces.

The \(\mathrm{Compat}\) fields of the original finite-level
residual data identify their local residual characters and their
\(\theta\) with the original ones. Thus \(u_N.\psi\) satisfies
the whole \(r.\mathrm{IsPsi}\). The \(U\)-congruences are
\((A.A_N).U_{\mathrm{res}}\). Finally all inner and outer
power-series variables are in the maximal ideal of the nested
\(\mathrm{Sinf}\), and their images under \(s_N\) remain there
by PRG2. This checks every field of the actual \texttt{tupN}.
No level is discarded and no eventual equality replaces these
exact equalities at each \(N\).
% END PRG19

% BEGIN PRG20
\paragraph{PRG20. The specialization tuple over the exact ordered splitting algebra.}
\label{claim:PRG20}
Write \(S_0=\mathrm{Spo}(A)\), \(\iota_B:B\to S_0\),
and \(u^0_{j,i}\) for its actual ordered roots. The following
is a tuple over \(S_0\):
\[
 \left((u.D\circ p)\otimes_B S_0,\
       \iota_B\circ u.\psi,\
       \iota_B(A.b.\lambda(X_{\mathrm{unit}})),\
       (u^0_{j,i}),\ (0)_{j,i}\right).                    \tag{PRG.16}
\]
By PRG18, \(S_0\) is in \(\mathcal C_O\). Continuity and
residue compatibility of \(\iota_B\) follow from PRG2; combined
with continuity of \(p\) they prove the determinant continuity
and residual identity. Applying \(\iota_B\) to the three
\(\mathrm{IsPsi}\) fields proves them for the new characters.
The value \(A.b.\lambda(X_{\mathrm{unit}})\) is in
\(\mathfrak m_B\) by PRG2 applied to its original power-series
map. Each root minus \(A.\mathrm{ut}_{j,i}\) is the image of
the corresponding variable in (PRG.13), so lies in
\(\mathfrak m_{S_0}\). The last coordinates are zero. These
checks prove exactly \texttt{tupPi}, with the original map from
\(B\) and the original roots. They require no choice of a
geometric point or of a root ordering other than the ordering
already in \(\mathrm{Sp}\).
% END PRG20

% BEGIN PRG21
\paragraph{PRG21. The universal section with the original weight coordinates.}
\label{claim:PRG21}
For any \(W\) as in PRG14, put
\[
 w_k=A.b.\mathrm{wt}.c_k\,
       W.t.\psi_{A.b.\mathrm{wt}.v_k,\,A.b.\mathrm{wt}.i_k}
                    (A.b.\mathrm{wt}.\gamma_k)-1.
 \tag{PRG.17}
\]
Its residue is zero: the character's residue field and the exact
\(A.b.\mathrm{wt\_res}\) identity give the product \(1\).
Thus \(w_k\in\mathfrak m_{W.R}\). The original inner ring is
\(L_0=O[[X_{\mathrm{unit}},(X_k)_{k<D}]]\), with its literal
\(\mathrm{Unit}\) coordinate. The original outer ring is
\(\mathrm{Sinf}=L_0[[(Y_{j,i})]]\). Apply PRG3 twice to obtain
a unique \(O\)-map
\[
 s_\infty:\mathrm{Sinf}\longrightarrow W.R,
 \quad CX_{\mathrm{unit}}\mapsto W.t.s,
 \quad CX_k\mapsto w_k,
 \quad Y_{j,i}\mapsto W.t.y_{j,i}.                         \tag{PRG.18}
\]
The needed maximal-ideal memberships are the tuple fields and
(PRG.17). Every map out of \(\mathrm{Sinf}\) to an object of
\(\mathcal C_O\) is determined by all these inner and outer
variable values: first restrict to \(L_0\) and use PRG3, then
apply it to the outer series. This proves the exact
\texttt{sinfExt} assertion and the full section fields of
\(\mathrm{RMaps}\), including the original weight factors
\(c_k,v_k,i_k,\gamma_k\).
% END PRG21

% BEGIN PRG22
\paragraph{PRG22. All finite-level maps and their identities on every series.}
\label{claim:PRG22}
Under adaptedness, apply \(W.\mathrm{univ}\) to the actual tuple
(PRG.15), for every \(N\). It gives an \(O\)-map
\(\rho_N:W.R\to R_N\) carrying the full determinant, all
characters, \(s\), every \(U_{j,i}\) and every \(y_{j,i}\)
to the five entries of (PRG.15). These are exactly
\(\rho_{\mathrm{maps}}\), with no alteration of the finite-level
rings or pair data. Moreover
\[
 \rho_N(s_\infty z)=s_N(z)\qquad
       (\text{every }N,\ \text{every }z\in\mathrm{Sinf}).   \tag{PRG.19}
\]
For the inner Unit variable and all outer variables this is their
defining tuple equality. For an inner weight variable, apply
\(\rho_N\) to (PRG.17). It fixes the coefficient \(c_k\),
carries the character value to the corresponding value of
\(u_N.\psi\), and therefore gives exactly the original
\((A.A_N).s_w\) formula for \(s_N(CX_k)\). The uniqueness in
PRG21 now proves equality of the two maps on all of
\(\mathrm{Sinf}\). In particular (PRG.19) is not limited to
polynomials or to a selected dense subalgebra. PRG2 supplies the
continuity needed to pass from variables to complete series.
% END PRG22

% BEGIN PRG23
\paragraph{PRG23. The exact specialization map and constant-coefficient identity.}
\label{claim:PRG23}
Apply \(W.\mathrm{univ}\) to (PRG.16). The resulting
\(\pi:W.R\to S_0\) carries all five tuple entries to their
displayed targets, which proves every field of
\(\pi_{\mathrm{maps}}\). If
\(\mathrm{cc}:L_0[[(Y_{j,i})]]\to L_0\) is constant
coefficient in the outer variables, then
\[
 \pi(s_\infty z)=\iota_B\bigl(A.b.\lambda(\mathrm{cc}(z))\bigr)
       \quad(\text{every }z\in\mathrm{Sinf}).              \tag{PRG.20}
\]
The maps agree on the inner Unit variable by the third entry of
(PRG.16), and both kill every outer variable by its last entry.
For the inner weight variable, the image of (PRG.17) is
\[
 \iota_B\bigl(c_k\,u.\psi_{v_k,i_k}(\gamma_k)-1\bigr)
       =\iota_B\bigl(A.b.\lambda(X_k)\bigr),
\]
by the original \(A.b.\mathrm{segmentWeightExponent}\) field.
Both maps are \(O\)-algebra maps to an object of
\(\mathcal C_O\), so PRG21 gives equality on every power
series. Notice that (PRG.20) uses the actual \(A.b.\lambda\)
and the actual inclusion \(B\to\mathrm{Spo}(A)\).
There is no replacement by a residue specialization and no new
ideal in place of \(J\).
% END PRG23

% BEGIN PRG24
\paragraph{PRG24. Complete \(\mathrm{RMaps}\) and the exact remaining patching contract.}
\label{claim:PRG24}
PRG21, PRG22 and PRG23 assemble, in the source's order, into
\[
 (s_\infty,\ s_{\mathrm{Unit}},\ s_w,\ s_y,\
       (\rho_N),\ \rho_{\mathrm{maps}},\
       (\rho_Ns_\infty=s_N),\ \pi,\ \pi_{\mathrm{maps}},\
       \pi s_\infty=\iota_B\lambda\mathrm{cc}).            \tag{PRG.21}
\]
Every field is an exact equality at its original quantifier scope.
Thus \texttt{rmapsEx} is proved for any original \(W\), assuming
precisely the original adaptedness condition. Together with
PRG13 this supplies actual \(W:\mathrm{RRObj}\) and
\(m:\mathrm{RMaps}(W)\) on the actual given limit \(\Lambda\)
and finite-level system \(A\). Their rings, maps, ordered roots,
weight factors, Unit coordinate and determinant laws all remain
those required by the source.

This closes the universal-ring and specialization-map part of
the limiting patching dependency. The further work is to construct
the limiting module from the original finite modules, prove the
six families of approximate annihilator relations and the full
support/completion assertion, and then apply them to the actual
quotient pair and point. Those are separate producers; no module
freeness, support dimension or final \(J\)-containment is inferred
merely from (PRG.21). In the selected arithmetic application
\(\ell=3\), \(n=4\), \(g=0\) and
\(D_0=4[L^+:\mathbf Q]+4q\) are retained. The preceding
construction is valid at these values without changing either
rank-four Link or the fixed complex embedding.
% END PRG24
